\section{Product charts and boxes}\label{sec:prelim}

Throughout the proof, $X,E_1,E_2$ satisfy the hypotheses of
Theorem~\ref{thm:main}. Write $\D=\{z\in\C:|z|<1\}$.
Fix a K\"ahler form $\omega$ on $X$ and its Hermitian metric
$g(a,b)=\omega(a,Jb)$ on the real tangent bundle, where $J$ is the
complex structure. We identify the real tangent bundle with the
holomorphic tangent bundle as complex vector bundles, so multiplication
by $i$ on real tangent vectors means application of $J$. The splitting
therefore has complex-linear projections $P_i$ on real tangent vectors.
Compactness gives
\begin{equation}\label{eq:projection-bounds}
\sup_{x\in X}\norm{P_i(x)}<\infty\qquad(i=1,2).
\end{equation}
No orthogonality of the two summands is assumed.

\begin{lemma}\label{lem:product-charts}
Every point of $X$ has a holomorphic coordinate neighborhood with
coordinates $z=(z_1,z_2)$ and range $\D^{r_1}\times\D^{r_2}$, in
which $E_i$ is tangent to the $z_i$-directions. Changes between these
charts are locally of the following form, with each $h_i$ a local
biholomorphism:
\begin{equation}\label{eq:separated-transition}
z'_1=h_1(z_1),\qquad z'_2=h_2(z_2).
\end{equation}
\end{lemma}

\begin{proof}
Holomorphic Frobenius gives a local submersion with kernel $E_i$ for
each $i$. To recall the local construction, choose a coordinate
projection that is nonsingular on the distribution and lift its
coordinate fields to a frame in the distribution. Their brackets
belong to the distribution by integrability and project to zero;
the projection is an isomorphism on the distribution, so the brackets
vanish. The commuting local holomorphic flows applied to a transverse
slice give foliation coordinates and the asserted submersion.

Combine the submersion with kernel $E_2$ and the submersion with kernel
$E_1$. Their combined differential is an isomorphism by complementarity.
The holomorphic inverse function theorem and restriction to a product
of polydisks give the desired coordinates. A chart change preserves
both distributions, so its wrong-block derivatives vanish. On small
product neighborhoods this is exactly
\eqref{eq:separated-transition}, and invertibility of the chart change
makes each block a local biholomorphism.
\end{proof}

We call these \emph{product charts}. A \emph{first plaque} is a slice
with $z_2$ fixed, and a \emph{second plaque} is a slice with $z_1$ fixed.
A continuous path is a first or second \emph{plaque path} if it locally
lies in plaques of the corresponding type. Transitions are required
to be separated only locally; no assertion is made about a single
formula on a disconnected chart overlap.

\begin{definition}\label{def:box}
Let $U\subset\C^a\times\C^b$ be open, with parameter blocks $(s,t)$.
A holomorphic map $f:U\to X$ is a \emph{holomorphic box} if
\[
\dd f(T\C^a\oplus0)\subset E_1,\qquad
\dd f(0\oplus T\C^b)\subset E_2.
\]
The conditions are imposed at every point of $U$. They do not require
$f$ to be an immersion.
\end{definition}

The definition imposes only tangency; the standing integrability
hypothesis supplies the product charts. In such a chart, a box is locally $(f_1(s),f_2(t))$. Conversely,
this local separation implies the two tangencies. The scalar case
$a=b=1$ will be used first. Its parameters $(s,t)$ retain their
prescribed directions even when we describe a domain using other
affine coordinates.

All uses of the identity theorem below are for holomorphic maps or
holomorphic bundle maps on connected domains. In particular, tangency
persists under holomorphic extension: the wrong-summand projections
of the differential vanish wherever the initial box is defined and
therefore throughout the connected extension domain.
